\documentclass[11pt,a4paper]{article}

\usepackage[utf8]{inputenc}
\usepackage[T1]{fontenc}
\usepackage[french]{babel}
\usepackage{lmodern}
\usepackage{microtype}
\usepackage{amsmath,amssymb,amsfonts,mathtools}
\usepackage{geometry}
\usepackage{hyperref}
\usepackage{xcolor}
\usepackage{enumitem}
\usepackage{booktabs}
\usepackage{array}
\usepackage{tcolorbox}

\geometry{margin=2.35cm}
\hypersetup{
  colorlinks=true,
  linkcolor=blue!55!black,
  urlcolor=blue!55!black,
  citecolor=blue!55!black,
  pdftitle={Conjecture UV/IR sur la constante de structure fine},
  pdfauthor={Notes de travail}
}

\setlength{\parskip}{0.6em}
\setlength{\parindent}{0pt}
\setlist[itemize]{leftmargin=1.6em}
\setlist[enumerate]{leftmargin=1.8em}

\newcommand{\lP}{\ell_{\mathrm P}}
\newcommand{\rhoP}{\rho_{\mathrm P}}
\newcommand{\RL}{R_{\Lambda}}
\newcommand{\SL}{S_{\Lambda}}
\newcommand{\Seuil}{S_0^{U(1)}}
\newcommand{\geff}{g_{\mathrm{eff}}^{U(1)}}
\newcommand{\gedge}{g_{\mathrm{edge}}^{U(1)}}
\newcommand{\apinv}{\alpha^{-1}}
\newcommand{\Cren}{\mathcal C_{\mathrm{ren}}}

\title{Conjecture UV/IR sur la constante de structure fine\\
\large Audit mathematique, test numerique et statut physique}
\author{Document de travail}
\date{\today}

\begin{document}
\maketitle

\begin{abstract}
Cette note formalise et verifie une relation candidate entre la constante de
structure fine, l'echelle de Planck et le rayon de de Sitter associe a la
constante cosmologique :
\[
\apinv=\ln\left(\frac{\RL}{10\pi\lP}\right).
\]
L'algebre interne de la relation est correcte et le test numerique avec
CODATA 2022 et Planck 2018 donne une proximite forte avec la valeur observee de
\(\alpha^{-1}\). Le statut scientifique reste cependant conditionnel : le
facteur \(10\pi\), ou de maniere equivalente
\(\geff=5/2\), doit etre derive independamment depuis la microphysique des
modes de bord Maxwell. Le resultat doit donc etre lu comme une conjecture
UV/IR precise, non comme une derivation etablie de \(\alpha\).
\end{abstract}

\tableofcontents
\newpage

\section{Resultat candidat}

On definit
\begin{equation}
\RL=\sqrt{\frac{3}{\Lambda}},
\qquad
\lP=\sqrt{\frac{\hbar G}{c^3}}.
\end{equation}

La relation candidate est
\begin{equation}
\boxed{\apinv=\ln\left(\frac{\RL}{10\pi\lP}\right).}
\label{eq:main_alpha}
\end{equation}

Elle est equivalente a
\begin{equation}
\boxed{\Lambda=\frac{3}{100\pi^2\lP^2}\,e^{-2/\alpha}.}
\label{eq:lambda_alpha}
\end{equation}

Cette formule convertit la petitesse de \(\Lambda\) en suppression
exponentielle controlee par \(\alpha\). Mais cette equivalence ne suffit pas a
constituer une preuve : le coefficient \(10\pi\) doit etre justifie sans
utiliser la valeur observee de \(\alpha\).

\section{Lien UV/IR de l'energie noire}

La densite associee a \(\Lambda\) est
\begin{equation}
\rho_\Lambda=\frac{\Lambda c^4}{8\pi G}
=\frac{3c^4}{8\pi G\RL^2}.
\end{equation}

Comme
\begin{equation}
\rhoP=\frac{c^7}{\hbar G^2},
\qquad
\rhoP\lP^2=\frac{c^4}{G},
\end{equation}
on obtient l'identite
\begin{equation}
\boxed{\rho_\Lambda=\frac{3}{8\pi}\rhoP
\left(\frac{\lP}{\RL}\right)^2.}
\label{eq:rho_lambda_uv_ir}
\end{equation}

Le meme rapport UV/IR apparait donc sous deux formes :
\begin{equation}
\frac{\rho_\Lambda}{\rhoP}\sim \left(\frac{\lP}{\RL}\right)^2,
\qquad
\apinv\sim \ln\left(\frac{\RL}{\lP}\right).
\end{equation}

\section{Forme entropique}

L'entropie d'horizon de de Sitter vaut
\begin{equation}
\SL=\frac{k_B A_\Lambda}{4\lP^2}
=\frac{k_B\pi\RL^2}{\lP^2}.
\end{equation}

La conjecture entropique est
\begin{equation}
\boxed{\apinv=\frac12\ln\left(\frac{\SL}{\Seuil}\right).}
\label{eq:alpha_entropy}
\end{equation}

Si
\begin{equation}
\Seuil=100\pi^3 k_B,
\label{eq:S0_value}
\end{equation}
alors
\begin{align}
\apinv
&=\frac12\ln\left(
\frac{k_B\pi\RL^2/\lP^2}{100\pi^3 k_B}
\right)\\
&=\frac12\ln\left(\frac{\RL^2}{100\pi^2\lP^2}\right)\\
&=\ln\left(\frac{\RL}{10\pi\lP}\right).
\end{align}

L'algebre est donc exacte. Le point non demontre est le seuil
\(\Seuil=100\pi^3k_B\).

\section{Flot unite}

On pose la capacite holographique
\begin{equation}
N(L)=\frac{S(L)}{k_B}=\frac{\pi L^2}{\lP^2}.
\end{equation}

Si
\begin{equation}
\alpha^{-1}(L)=\frac12\ln\left(\frac{N(L)}{N_0}\right),
\end{equation}
alors
\begin{equation}
\alpha^{-1}(L)
=\frac12\ln\left(\frac{\pi L^2}{\lP^2N_0}\right),
\end{equation}
et donc
\begin{equation}
\boxed{\frac{d\alpha^{-1}}{d\ln L}=1.}
\label{eq:flow_unit}
\end{equation}

Ce flot unite n'est pas ajuste. Il decoule de la loi d'aire
\(S\propto L^2\) et du facteur \(1/2\) dans la definition entropique.

\section{Reduction du facteur \texorpdfstring{\(10\pi\)}{10 pi}}

On parametre le seuil par
\begin{equation}
\frac{\Seuil}{k_B}=4\pi\left(2\pi\geff\right)^2.
\label{eq:S0_geff}
\end{equation}

En imposant \(\Seuil/k_B=100\pi^3\), on obtient
\begin{align}
4\pi(2\pi \geff)^2&=100\pi^3,\\
16\pi^3(\geff)^2&=100\pi^3,\\
(\geff)^2&=\frac{25}{4}.
\end{align}

Donc
\begin{equation}
\boxed{\geff=\frac52.}
\label{eq:geff_52}
\end{equation}

Le facteur macroscopique devient
\begin{equation}
4\pi\geff=10\pi.
\end{equation}

Le probleme numerique du \(10\pi\) est donc reduit au probleme
microphysique de deriver \(\geff=5/2\).

\section{Interpretation \texorpdfstring{\(2+1/2\)}{2+1/2}}

Le photon libre en quatre dimensions possede deux polarisations physiques :
\begin{equation}
N_{\mathrm{pol}}=2.
\end{equation}

Dans une region avec bord, la contrainte de Gauss empeche une factorisation
naive de l'espace de Hilbert. Le flux electrique normal au bord devient une
donnee physique de bord. De maniere schematique, la forme symplectique de bord
contient un couple conjugue
\begin{equation}
\Omega_{\mathrm{edge}}
=\int_{\partial\Omega}\delta E_\perp\wedge\delta\varphi.
\end{equation}

Une integration gaussienne reelle produit une demi-puissance de determinant :
\begin{equation}
\int d^n x\,\exp\left(-\frac12x^TKx\right)
=(2\pi)^{n/2}(\det K)^{-1/2}.
\end{equation}

Cela motive le comptage effectif
\begin{equation}
\boxed{\geff=2+\frac12=\frac52.}
\label{eq:2_plus_half}
\end{equation}

Cette ligne est plausible, mais elle n'est pas encore une demonstration. Les
modes de bord Maxwell existent dans les calculs d'entropie, mais l'association
entre une demi-puissance de determinant et une capacite positive \(+1/2\) est
une regle supplementaire.

\section{Capacite holographique renormalisee}

Sur une surface \(\Sigma\), le facteur de bord peut prendre la forme
\begin{equation}
Z_{\mathrm{edge}}\propto
\left[\det{}'(-\Delta_\Sigma)\right]^{q},
\qquad q=\pm\frac12.
\end{equation}

La puissance demi-entiere est robuste, mais le signe, les modes zero, la mesure
fonctionnelle et les contretermes dependent du schema. On propose donc une
capacite renormalisee
\begin{equation}
\boxed{\Cren[(\det K)^q]=|q|.}
\label{eq:Cren}
\end{equation}

Cette definition donne
\begin{equation}
\Cren[Z_{\mathrm{edge}}]=\frac12.
\end{equation}

Le verrou theorique est precis : il faut montrer que cette capacite
renormalisee est la bonne quantite physique entrant dans la relation
entropique de \(\alpha\).

\section{Theoreme conditionnel}

\textbf{Theoreme conditionnel.} Supposons que :
\begin{enumerate}
\item la capacite holographique d'un horizon de rayon \(L\) soit
\[
N(L)=\frac{\pi L^2}{\lP^2};
\]
\item l'inverse du couplage electromagnetique soit la distance logarithmique
entropique
\[
\alpha^{-1}(L)=\frac12\ln\left(\frac{N(L)}{N_0^{U(1)}}\right);
\]
\item le seuil \(U(1)\) soit
\[
N_0^{U(1)}=4\pi\left(2\pi\geff\right)^2;
\]
\item la microphysique de bord donne \(\geff=5/2\).
\end{enumerate}

Alors
\begin{equation}
N_0^{U(1)}=100\pi^3,
\end{equation}
et
\begin{equation}
\boxed{\alpha^{-1}(L)=\ln\left(\frac{L}{10\pi\lP}\right).}
\end{equation}

En choisissant \(L=\RL\), on obtient \eqref{eq:main_alpha}.

\section{Test numerique}

Avec les constantes CODATA 2022 et les valeurs centrales Planck 2018
\[
H_0=67.4\,{\rm km\,s^{-1}\,Mpc^{-1}},
\qquad
\Omega_\Lambda=0.685,
\]
on prend
\begin{equation}
\Lambda_{\rm obs}=\frac{3\Omega_\Lambda H_0^2}{c^2},
\qquad
\RL=\sqrt{\frac{3}{\Lambda_{\rm obs}}}.
\end{equation}

Le script du dossier donne :
\begin{center}
\begin{tabular}{lr}
\toprule
Quantite & Valeur \\
\midrule
\(\Lambda_{\rm obs}\) & \(1.090910502838\times10^{-52}\,{\rm m^{-2}}\)\\
\(\RL\) & \(1.658311322028\times10^{26}\,{\rm m}\)\\
\(\alpha^{-1}_{\rm CODATA}\) & \(137.035999177\)\\
\(\alpha^{-1}_{\rm pred}\) & \(137.036063742708\)\\
\(\Delta\alpha^{-1}\) & \(6.456570762\times10^{-5}\)\\
\(\Lambda(\alpha)/\Lambda_{\rm obs}\) & \(1.000129139753\)\\
\(g_{\rm eff,fit}\) & \(2.500161419480\)\\
\bottomrule
\end{tabular}
\end{center}

La relation est donc numeriquement tres proche. Cependant, l'incertitude
experimentale sur \(\alpha^{-1}\) est beaucoup plus petite que cet ecart. La
proximite doit etre lue comme un indice de structure, pas comme une validation
experimentale.

\section{Conclusion}

Les identites mathematiques de la note sont correctes :
\begin{equation}
\frac12\ln\left(\frac{\SL}{100\pi^3 k_B}\right)
=\ln\left(\frac{\RL}{10\pi\lP}\right),
\qquad
\frac{d\alpha^{-1}}{d\ln L}=1.
\end{equation}

Le test public donne une coincidence forte :
\begin{equation}
\alpha^{-1}_{\rm pred}=137.036063742708,
\qquad
\alpha^{-1}_{\rm obs}=137.035999177.
\end{equation}

Le point decisif reste ouvert :
\begin{equation}
\boxed{\text{deriver independamment }\geff=2+\frac12.}
\end{equation}

Si ce calcul de bord Maxwell est etabli dans une capacite holographique
renormalisee, la relation \(\alpha\)--\(\Lambda\) deviendra une prediction
structurelle du cadre UV/IR. En l'etat, elle constitue une conjecture precise,
testable et mathematiquement coherente, mais pas encore une publication
definitive.

\begin{thebibliography}{9}

\bibitem{CODATA2022}
NIST/CODATA, \emph{CODATA recommended values of the fundamental physical
constants: 2022}.

\bibitem{Planck2018}
Planck Collaboration, \emph{Planck 2018 results. VI. Cosmological parameters},
arXiv:1807.06209.

\bibitem{DonnellyWall2015}
W. Donnelly and A. C. Wall, \emph{Geometric entropy and edge modes of the
electromagnetic field}, arXiv:1506.05792.

\bibitem{DonnellyWallPRL2015}
W. Donnelly and A. C. Wall, \emph{Entanglement entropy of electromagnetic edge
modes}, Phys. Rev. Lett. 114, 111603 (2015).

\end{thebibliography}

\end{document}
